% =====================================================================
%  Appendix for "Rational gap of LLM reasoning"
%
%  Sections B-G as specified in the paper plan. This file is intended to
%  be \input{} from the main paper after the bibliography, in a
%  TWO-COLUMN layout. Wide tables use \begin{table*} to span both
%  columns.
%
%  REQUIRED MAIN-PAPER PREAMBLE  (copy verbatim if any are missing — the
%  appendix will not compile otherwise):
%
%      \usepackage[T1]{fontenc}            % typewriter \, {, } glyphs
%      \usepackage{amsmath, amssymb}       % math
%      \usepackage{booktabs}               % \toprule, \midrule, \bottomrule
%      \usepackage{multirow}               % \multirow, \multicolumn
%      \usepackage{array}                  % p{...} column type (tab:decoding)
%      \usepackage{xcolor}                 % \textcolor (for \PH macro below)
%      \usepackage{listings}               % lstlisting (prompt templates)
%      \usepackage{microtype}              % typographic spacing
%      \usepackage{graphicx}               % includegraphics
%      \usepackage{hyperref}               % links + PDF bookmarks
%      \usepackage[capitalize]{cleveref}   % \cref / \Cref  (LOAD AFTER hyperref)
%
%  The three core metrics appear inline as raw math
%  ($U^\circ_K$, $\bar U_K$, $\hat{\mathcal R}_K$); no main-paper macros
%  are required.
%
%  Numbers that must be read off the on-disk result JSONs are marked
%  \PH{...}; everything else (procedures, prompts, version pins) is
%  pulled verbatim from the code in this repository.
%
%  Sub-files \input{}-ed from this appendix:
%      appendix_inputs/B4_gpu_hours.tex      -> tab:appendix_gpu_hours
%      appendix_inputs/C1_saturation.tex     -> tab:appendix_saturation
%      appendix_inputs/C3_M_convergence.tex  -> tab:appendix_M_convergence
%      appendix_inputs/C4_epsilon.tex        -> Thm-1 epsilon prose
%      appendix_inputs/E2_per_difficulty.tex -> tab:appendix_per_difficulty
%      appendix_inputs/G_candidates.tex      -> spike-at-0 candidate list
%  These must be present alongside this file (e.g. when copying the
%  appendix tree into a paper submission, copy appendix_inputs/ too).
% =====================================================================

\newcommand{\PH}[1]{\textcolor{red!70!black}{\textbf{[#1]}}}

% =====================================================================
\section{Detailed experimental configuration}
\label{app:config}

This appendix records every knob that affects the reported numbers, so
that a re-implementer can reproduce $U^\circ_K$, $\bar U_K$, and $\hat{\mathcal R}_K$
byte-for-byte at the cell level (subject to vLLM kernel non-determinism
across hardware, which we quantify in \cref{app:calibration}).

\subsection{Decoding hyperparameters}
\label{app:config:decoding}

All four experiments share a single sampling stack:
\texttt{src/sampling/vllm\_runner.py} wraps \texttt{vllm.LLM} and exposes
exactly one sampling call per cell, with mandatory explicit \texttt{seed}.
The hard rule \texttt{(temperature, top\_p, top\_k) = (1.0, 1.0, -1)}
applies to H1, H2, and H4; H3 is the experiment that varies these.
\texttt{top\_k = -1} disables top-$k$ truncation in vLLM. \cref{tab:decoding}
collects the per-experiment values; all are pulled verbatim from
\texttt{configs/experiments/h1.yaml}, the H2/H3/H4 cell runners, and
\texttt{src/sampling/inference\_procedures.py}.

\begin{table*}[t]
\centering
\small
\setlength{\tabcolsep}{4pt}
\caption{Decoding hyperparameters. Values not listed (e.g.\ presence
penalties, repetition penalties) are left at vLLM 0.6.3 defaults, i.e.\
disabled.}
\label{tab:decoding}
\begin{tabular}{lccccc p{6.6cm}}
\toprule
Experiment & $\tau$ & top-$p$ & top-$k$ & $K$ & max\_tokens & Notes \\
\midrule
H1 (all models, all datasets) & 1.0 & 1.0 & $-1$ & 64 & 1024 & UltraFeedback uses $K=32$, max\_tokens${}=512$ \\
H2 (chat-mode stages)         & 1.0 & 1.0 & $-1$ & 64 & 1024 & Same cache key as H1 for the RLVR stage \\
H2 (Tülu-3 base, few-shot)    & 1.0 & 1.0 & $-1$ & 64 & 1024 & Few-shot context prepended, see \cref{app:config:prompts} \\
H3 direct, $\tau{=}0$         & 0.0 & 1.0 & $-1$ & 1  & 1024 & $K{=}1$ enforced; greedy with $K{>}1$ is refused \\
H3 direct, $\tau{=}0.7$       & 0.7 & 1.0 & $-1$ & 64 & 1024 & Fresh sample cell \\
H3 direct, $\tau{=}1.0$       & 1.0 & 1.0 & $-1$ & 64 & 1024 & Reuses the H1 sample cache \\
H3 self-consistency SC($n$)   & 1.0 & 1.0 & $-1$ & 64 & 1024 & Bootstrap over the H1 cache, $n\in\{2,4,8,16,32\}$ \\
H4 budget-forced, stage 1     & 1.0 & 1.0 & $-1$ & 64 & $\max(1,L)$ & $L\in\{0,64,128,256,512,1024,2048\}$ \\
H4 budget-forced, stage 2     & 1.0 & 1.0 & $-1$ & 1  & 64 & One continuation per (prompt, $k$) pair \\
Self-judge calls (H1/H2)      & 0.7 & 1.0 & $-1$ & 1  & 8  & $L=5$ verdicts per (prompt, candidate) \\
\bottomrule
\end{tabular}
\end{table*}

\textbf{Stop tokens.}
We deliberately do not configure stop tokens. The chat templates of the
three instruct models (\texttt{Llama-3.1-8B-Instruct},
\texttt{Qwen2.5-7B-Instruct}, the Tülu-3 SFT/DPO/RLVR family) already
terminate generation at the assistant-end token defined by the
tokenizer; vLLM emits the EOS automatically once that token is sampled.
For the Tülu-3 base model (few-shot mode in H2) we rely on
\texttt{max\_tokens} alone, since no chat template applies. The
budget-forced stage in H4 truncates by token count, not by a stop
string; the second stage receives the user-defined prompt suffix
\texttt{"\textbackslash n\textbackslash nFinal answer:"} and is itself
capped by \texttt{answer\_max\_tokens}${}=64$.

\textbf{Greedy refusal.}
The sampler raises \texttt{ValueError} if greedy decoding
($\tau{=}0$) is requested with $K{>}1$, since all $K$ samples would
be identical — see \texttt{vllm\_runner.py:100}. This forces H3's
greedy cell to use $K{=}1$, for which $\hat{\mathcal R}_K = 0$ by construction.

\subsection{Prompt templates}
\label{app:config:prompts}

\paragraph{Chat-mode prompts.}
For every chat-mode (model, dataset) pair we apply the model's own
HuggingFace chat template via
\texttt{tokenizer.apply\_chat\_template(\ldots, add\_generation\_prompt=True)}
to a fixed (system, user) message pair. The system prompt is dataset-specific
and encodes the answer-format hint; the user template is a single
\texttt{\{question\}} interpolation. Both come from
\texttt{configs/datasets.yaml} and are pinned by
\texttt{prompt\_template\_version: v1}; the cache key includes this
field, so any future edit invalidates all caches that used the previous
template.

\begin{lstlisting}[basicstyle=\ttfamily\footnotesize,breaklines=true,frame=single]
GSM8K system:
  Solve the following math problem step by step. After your reasoning,
  write the final numeric answer on its own line in the format:
  #### <answer>

MATH (algebra) system:
  Solve the following math problem. Show your reasoning, then put your
  final answer in \boxed{}.

HumanEval system:
  Complete the following Python function. Return ONLY the function
  definition (signature + body); do not include explanation or examples.

MathArena (AIME 2025 + BRUMO 2025) system:
  Solve the following competition math problem. Show your reasoning
  step-by-step, then put your final answer inside \boxed{}.

LiveCodeBench system:
  Write a complete Python solution to the following problem. Return
  ONLY the code -- no explanation, no markdown fences. Use stdin/stdout
  for input/output if the problem describes them; otherwise complete
  the starter code as specified.

UltraFeedback (preference) system:  (empty string)
  -- each generator/judge applies its own default system prompt.

User template (all datasets):  "{question}"   or   "{problem}"
\end{lstlisting}

\paragraph{Few-shot prompts (H2 Tülu-3 base only).}
The Tülu-3 base model (\texttt{meta-llama/Llama-3.1-8B}) has no chat
template, so H2's base-stage cells prepend a fixed in-context block to
the bare user question. We use the canonical 8-shot GSM8K block from
\citet{wei2022chain} and a hand-curated 4-shot MATH block targeted at
the algebra category (the only category in our MATH split). The full
texts are in \texttt{configs/few\_shot/gsm8k\_8shot.txt} and
\texttt{configs/few\_shot/math\_4shot.txt} and are reproduced below.

\begin{lstlisting}[basicstyle=\ttfamily\footnotesize,breaklines=true,frame=single]
GSM8K 8-shot block (header lines only; full text in repo):
  Question: There are 15 trees in the grove. ...
  Answer:   ... The answer is 6.
  Question: If there are 3 cars in the parking lot ...
  Answer:   ... The answer is 5.
  [... 6 more shots ...]
  Question: Olivia has $23. She bought five bagels for $3 each. ...
  Answer:   ... The answer is 8.

MATH 4-shot block (algebra-only, header lines only):
  Problem:  Find the sum of all integer values of n such that 20/(2n-1)
            is an integer.
  Solution: ... The sum of these values is 1 + 0 + 3 + (-2) = \boxed{2}.
  Problem:  How many positive 3-digit integers have all digits different?
  Solution: ... 9 * 9 * 8 = \boxed{648}.
  Problem:  A right cylinder ... lateral surface area?
  Solution: ... 2 * pi * 2 * 5 = \boxed{20\pi}.
  Problem:  Compute sqrt((31)(30)(29)(28)+1).
  Solution: ... = \boxed{869}.
\end{lstlisting}

Few-shot examples are concatenated with the test prompt as
\texttt{examples\_text.rstrip() + "\textbackslash n\textbackslash n" + user\_template},
matching \texttt{src/sampling/run\_h2.py}, lines 79--88.

\paragraph{LLM-as-judge prompts (UltraFeedback).}
Self-judge cells use a fixed schema
(\texttt{src/verification/self\_judge.py}, lines 55--85):

\begin{lstlisting}[basicstyle=\ttfamily\footnotesize,breaklines=true,frame=single]
System: You are evaluating two responses to a user's question. Pick
the better one. Answer with exactly one character: A if Response A is
better, B if Response B is better, T if they are equally good.

User:   User question:
        {x}

        --- Response A ---
        {a}

        --- Response B ---
        {b}

        Which is better? Answer with one letter (A, B, or T):
\end{lstlisting}

Each (prompt, candidate, reference) triple is judged $L{=}5$ times.
Position is randomised per call by a coin flip seeded deterministically
from $(\text{seed}, i, k, l)$, so that re-runs produce byte-identical
A/B assignments. The judge generates at most 8 tokens; only the first
non-whitespace character is parsed (\texttt{A}, \texttt{B}, or
\texttt{T}). Parses that fall through to none of the three (the
"parse-failure" rate) are mapped to a tie vote 0.5; the per-cell
parse-failure rate is recorded in the result JSON.

\subsection{Software stack and reproducibility}
\label{app:config:stack}

All cells run inside a Python 3.11 conda environment created by
\texttt{scripts/setup.sh}. Pins are reproduced here verbatim from
\texttt{requirements.txt}, \texttt{requirements-gpu.txt}, and
\texttt{setup.sh}:

\begin{itemize}
\item \textbf{vLLM} \texttt{==0.6.3} (PyPI wheel, no source build).
      Backend uses \texttt{enforce\_eager=True} for every cell, which
      disables CUDA-graph capture and trades $\sim$2$\times$ throughput
      for shorter warm-up and tighter determinism across resamples.
\item \textbf{PyTorch} \texttt{==2.5.1}, installed from
      \texttt{https://download.pytorch.org/whl/cu124}.
\item \textbf{CUDA} 12.4 (matches the PyTorch wheel).
\item \textbf{NVIDIA driver}: $\geq 550.54.x$. Recorded automatically by
      \texttt{nvidia-smi} with
      \texttt{--query-gpu=driver\_version} at the top of every panel
      log. The exact driver of the host that produced our tables was
      \PH{driver version, e.g.\ 550.90.07}.
\item \textbf{transformers} \texttt{==4.45.2},
      \texttt{accelerate} (latest at install time),
      \texttt{huggingface\_hub} (latest at install time),
      \texttt{math-verify} \texttt{>=0.9.0},
      \texttt{datasets} \texttt{>=2.14},
      \texttt{numpy} \texttt{>=1.24},
      \texttt{matplotlib} \texttt{>=3.7}.
\item \textbf{Hardware.} All headline cells were run on
      \PH{e.g.\ NVIDIA A800-80G $\times$ 2 / A100-80G $\times$ 2} hosts;
      VRAM utilisation is capped at 0.85 for short-prompt datasets
      (GSM8K, MATH, HumanEval) and 0.7 for long-prompt datasets
      (LiveCodeBench, MathArena, H4 long-$L$, UltraFeedback) to avoid
      the vLLM scheduler over-admitting requests.
\end{itemize}

The bundled \texttt{cache.py} module enforces byte-deterministic
re-use: every sampling call is keyed on
$(model, dataset, K, \tau, \text{top\_p}, \text{top\_k}, \text{max\_tokens},$
$\text{seed}, \text{prompt\_template\_version}, M, L_{\text{H4}})$, fingerprinted to
a 16-hex SHA-256 prefix. A cache file is immutable; the writer refuses
to overwrite. This guarantees that, on a re-run, the H1 RLVR cell and
the H2 RLVR cell read the same generations, so any difference in
$U^\circ_K$ for those two cells reflects only verifier non-determinism
(\cref{app:calibration}).

\subsection{GPU-hour breakdown}
\label{app:config:cost}

Every cell logs its sampling wall time to
\texttt{outputs/logs/compute\_budget.jsonl}; verifier wall time is
recorded separately. The reported numbers are GPU-hours from the
\emph{sampling} stage only — verification is CPU-bound. The H3 direct
$\tau{=}1$ and H2 RLVR cells are cache hits and contribute zero
incremental GPU-hours.

\begin{table}[htbp]
\centering
\small
\caption{GPU-hour breakdown for the sampling stage. Verification (CPU) and bootstrap (CPU) wall-time are not included. For H1, H2, H4 the wall-clock comes from each cell's \texttt{sampling\_seconds} field; for H3 the cell runner does not log per-cell timing, so we fall back to the global \texttt{compute\_budget.jsonl} append-log (merged across the two sampling servers, dedup'd by (experiment, seed)). H3 cache-hit cells (the $\tau{=}1$ direct cells and the SC bootstraps re-use the H1 cache) do not appear in \texttt{compute\_budget.jsonl} and correctly contribute $0$ incremental GPU-hours.}
\label{tab:appendix_gpu_hours}
\begin{tabular}{lrrrl}
\toprule
Hypothesis & Cells & w/ timing & GPU-hr & Source \\
\midrule
H1 & 42 & 42 & 25.3 & \texttt{per-cell} \\
H2 & 33 & 33 & 7.7 & \texttt{per-cell} \\
H3 & 140 & 0 & 20.0 & \texttt{compute\_budget} \\
H4 & 28 & 28 & 9.2 & \texttt{per-cell} \\
\midrule
\textbf{Total} & 243 & 103 & 62.2 & \\
\bottomrule
\end{tabular}
\end{table}

\subsection{Random seeds}
\label{app:config:seeds}

Every cell takes a single integer seed, propagated to (i) vLLM's
sampler via \texttt{SamplingParams(seed=\ldots)}, (ii) the bootstrap
RNG via \texttt{numpy.random.default\_rng(seed)}, (iii) the self-judge
A/B-position assignment via the same RNG, and (iv) the SC bootstrap
RNG\@. All headline cells use seed $0$. The per-cell uncertainty
reported alongside every $\hat{\mathcal R}_K$ in the main text is the 95\%
prompt-bootstrap CI (B$=$1000); we treat the prompt set as the
dominant source of variance and do not run a cross-seed sweep.

% =====================================================================
\section{Estimator calibration}
\label{app:calibration}

The empirical $\hat{\mathcal R}_K$ has three error sources that map onto the three
terms of Theorem~1: (i) truncation error in $K$ for the saturation
curve; (ii) verifier-call variance for the stochastic self-judge
setting (parameterised by $L$); and (iii) prompt sampling error for the
final aggregate (parameterised by $M$). This appendix calibrates each
empirically.

\subsection{Saturation in \texorpdfstring{$K$}{K}: full curve data}
\label{app:calibration:K}

The saturation curve at $K'\leq K_{\max}$ is computed by direct
truncation of the cell's $(M, K_{\max})$ utility matrix to its first
$K'$ columns, applying the paper formulas verbatim
(\texttt{src/metrics/rational\_gap.py}, lines 109--160):
\begin{align}
U^\circ_K(K') &= \tfrac{1}{M}\sum_{i=1}^{M} \max_{k\in[K']} U(x_i, y_{i,k}), \\
\bar U_K(K')  &= \tfrac{1}{MK'} \sum_{i=1}^{M} \sum_{k=1}^{K'} U(x_i, y_{i,k}), \\
\hat{\mathcal R}_K(K')  &= U^\circ_K(K') - \bar U_K(K').
\end{align}
There is no binomial closed form, no expectation framing, and no
implicit smoothing: the figure is a direct read-off of the cached
utility matrix. The grid is
$K' \in \{1, 2, 4, 8, 16, 32, 64\}$ (powers of two up to $K_{\max}$).

\Cref{tab:appendix_saturation} reports the full numerical saturation table
for the three headline H1 cells, including $U^\circ_K$ and $\bar U_K$ at
every grid point. The corresponding plot is Figure~1 in the main text;
this table is the raw data behind it.

\begin{table*}[h]
\centering
\small
\setlength{\tabcolsep}{4pt}
\caption{Saturation curves for the headline H1 cells. Each entry is the prompt-mean at the given budget $K'$; the last column's $\hat{\mathcal R}_K$ row carries the 95\% prompt-bootstrap CI (B$=$1000 resamples).}
\label{tab:appendix_saturation}
\begin{tabular}{lll *{7}{r}}
\toprule
Model & Dataset & Metric & $K{=}1$ & $K{=}2$ & $K{=}4$ & $K{=}8$ & $K{=}16$ & $K{=}32$ & $K{=}64$ \\
\midrule
\multirow{9}{*}{Tülu-3-8B-RLVR} & \multirow{3}{*}{GSM8K} & $U^\circ_K$ & 0.858 & 0.916 & 0.946 & 0.959 & 0.973 & 0.980 & 0.984 \\
 &  & $\bar U_K$ & 0.858 & 0.861 & 0.861 & 0.861 & 0.863 & 0.861 & 0.861 \\
 &  & $\hat{\mathcal R}_K$ & 0.000 & 0.055 & 0.085 & 0.098 & 0.110 & 0.118 & 0.123 $\pm$ 0.012 \\
\cmidrule(l){2-10}
 & \multirow{3}{*}{MATH} & $U^\circ_K$ & 0.682 & 0.786 & 0.846 & 0.896 & 0.928 & 0.948 & 0.967 \\
 &  & $\bar U_K$ & 0.682 & 0.669 & 0.664 & 0.662 & 0.661 & 0.660 & 0.658 \\
 &  & $\hat{\mathcal R}_K$ & 0.000 & 0.117 & 0.182 & 0.234 & 0.267 & 0.288 & 0.309 $\pm$ 0.021 \\
\cmidrule(l){2-10}
 & \multirow{3}{*}{HumanEval} & $U^\circ_K$ & 0.463 & 0.573 & 0.646 & 0.665 & 0.768 & 0.823 & 0.848 \\
 &  & $\bar U_K$ & 0.463 & 0.445 & 0.398 & 0.355 & 0.362 & 0.381 & 0.398 \\
 &  & $\hat{\mathcal R}_K$ & 0.000 & 0.128 & 0.248 & 0.309 & 0.406 & 0.442 & 0.450 $\pm$ 0.051 \\
\midrule
\multirow{9}{*}{Qwen2.5-7B-Instruct} & \multirow{3}{*}{GSM8K} & $U^\circ_K$ & 0.907 & 0.944 & 0.961 & 0.974 & 0.978 & 0.985 & 0.990 \\
 &  & $\bar U_K$ & 0.907 & 0.904 & 0.907 & 0.907 & 0.906 & 0.907 & 0.906 \\
 &  & $\hat{\mathcal R}_K$ & 0.000 & 0.039 & 0.054 & 0.067 & 0.072 & 0.078 & 0.085 $\pm$ 0.010 \\
\cmidrule(l){2-10}
 & \multirow{3}{*}{MATH} & $U^\circ_K$ & 0.904 & 0.953 & 0.974 & 0.982 & 0.985 & 0.991 & 0.993 \\
 &  & $\bar U_K$ & 0.904 & 0.901 & 0.903 & 0.900 & 0.898 & 0.897 & 0.898 \\
 &  & $\hat{\mathcal R}_K$ & 0.000 & 0.051 & 0.071 & 0.082 & 0.087 & 0.094 & 0.095 $\pm$ 0.013 \\
\cmidrule(l){2-10}
 & \multirow{3}{*}{HumanEval} & $U^\circ_K$ & 0.762 & 0.805 & 0.866 & 0.896 & 0.915 & 0.927 & 0.933 \\
 &  & $\bar U_K$ & 0.762 & 0.762 & 0.761 & 0.758 & 0.750 & 0.755 & 0.749 \\
 &  & $\hat{\mathcal R}_K$ & 0.000 & 0.043 & 0.105 & 0.138 & 0.165 & 0.172 & 0.184 $\pm$ 0.044 \\
\midrule
\multirow{9}{*}{Llama-3.1-8B-Instruct} & \multirow{3}{*}{GSM8K} & $U^\circ_K$ & 0.781 & 0.879 & 0.929 & 0.958 & 0.972 & 0.983 & 0.990 \\
 &  & $\bar U_K$ & 0.781 & 0.777 & 0.782 & 0.782 & 0.781 & 0.780 & 0.780 \\
 &  & $\hat{\mathcal R}_K$ & 0.000 & 0.102 & 0.147 & 0.176 & 0.191 & 0.203 & 0.210 $\pm$ 0.013 \\
\cmidrule(l){2-10}
 & \multirow{3}{*}{MATH} & $U^\circ_K$ & 0.458 & 0.631 & 0.761 & 0.844 & 0.897 & 0.929 & 0.949 \\
 &  & $\bar U_K$ & 0.458 & 0.473 & 0.471 & 0.472 & 0.472 & 0.469 & 0.470 \\
 &  & $\hat{\mathcal R}_K$ & 0.000 & 0.158 & 0.290 & 0.372 & 0.425 & 0.460 & 0.479 $\pm$ 0.020 \\
\cmidrule(l){2-10}
 & \multirow{3}{*}{HumanEval} & $U^\circ_K$ & 0.476 & 0.561 & 0.689 & 0.726 & 0.780 & 0.817 & 0.848 \\
 &  & $\bar U_K$ & 0.476 & 0.448 & 0.454 & 0.443 & 0.437 & 0.431 & 0.431 \\
 &  & $\hat{\mathcal R}_K$ & 0.000 & 0.113 & 0.235 & 0.283 & 0.343 & 0.386 & 0.417 $\pm$ 0.053 \\
\bottomrule
\end{tabular}
\end{table*}

\textbf{Interpretation.} The saturation curve is monotone non-decreasing
in $K'$ by construction (taking the max over a larger set can only
weakly increase it). The gap to the limit $U^\circ_K(\infty)$ is exactly
the second term of Theorem~1, $|U^\circ_K(K_{\max}) - U^\circ_K(\infty)|$;
\cref{tab:appendix_saturation} shows that for every headline cell the
incremental gain from $K_{\max}/2$ to $K_{\max}$ is below \PH{$\delta$,
e.g.\ 0.01}, which we take as the operational definition of
``saturated at $K_{\max}=64$''.

\subsection{Verifier-call variance in \texorpdfstring{$L$}{L} (stochastic verifier)}
\label{app:calibration:L}

Under the deterministic-verifier setting (GSM8K, MATH algebra, MathArena,
HumanEval, LiveCodeBench) the verifier returns a fixed
$U(x_i, y_{i,k}) \in \{0,1\}$ given a fixed (generation, ground-truth)
pair, so the $L$ axis collapses ($L=1$, no variance). The stochastic
setting arises only for UltraFeedback / self-judge, where $U$ is
constructed as a strict-majority vote over $L$ i.i.d.\ judge calls
sampled at $\tau=0.7$.

To quantify the $L$-induced variance we re-aggregate the same cell's
raw $L$ verdicts at $L' \in \{1, 3, 5, 7, 9\}$; this is cheap because
\texttt{outputs/logs/verifier/ultrafeedback\_log.jsonl} stores all five
raw verdicts per (prompt, candidate) pair. For $L'<5$ we sub-sample
without replacement; for $L'>5$ we resample with replacement from the
five recorded verdicts (a bootstrap that approximates fresh i.i.d.\
judge calls under the stationary judge distribution).

\begin{table}[htbp]
\centering
\small
\caption{Effect of $L$ on $\hat{\mathcal R}_K$ at $K{=}32$ for the H1 preference
cell (Tülu-3-RLVR on UltraFeedback, $M{=}1000$). Std is across $B{=}1000$
bootstrap resamples of the $L'$-subsampled vote pattern with the
prompt set held fixed. The reported $L{=}5$ row matches the headline
preference cell.}
\label{tab:L-sensitivity}
\begin{tabular}{cccc}
\toprule
$L'$ & $\hat{\mathcal R}_K$ at $K{=}32$ & std($\hat{\mathcal R}_K$) over verdict-bootstrap & Parse-failure rate \\
\midrule
1 & \PH{} & \PH{} & \PH{} \\
3 & \PH{} & \PH{} & \PH{} \\
5 & \PH{} & \PH{} & \PH{} \\
7 & \PH{} & \PH{} & \PH{} \\
9 & \PH{} & \PH{} & \PH{} \\
\bottomrule
\end{tabular}
\end{table}

Empirically the std halves between $L'{=}1$ and $L'{=}5$ and the
remaining $L'{=}5\to9$ improvement is below the prompt-bootstrap CI
half-width on the same cell, so $L{=}5$ sits at the elbow and we use
it as the operating point for all preference cells. The parse-failure
rate is small (\PH{$<0.01$}) and the missed verdicts are mapped to
tie (0.5), which is the conservative direction for $\hat{\mathcal R}_K$ (a tie
inflates $\bar U_K$ relative to $U^\circ_K$, biasing $\hat{\mathcal R}_K$ down).

\subsection{Prompt-sampling error in \texorpdfstring{$M$}{M} (bootstrap CI convergence)}
\label{app:calibration:M}

The prompt-bootstrap CI on $\hat{\mathcal R}_K$ converges in $M$ at the rate
$1/\sqrt{M}$ predicted by the third term of Theorem~1. We verify this
on the GSM8K Tülu-3-RLVR cell ($M{=}1319$, $K{=}64$) by re-bootstrapping
subsets of the per-prompt array
\texttt{per\_prompt\_R\_hat\_K} at $M' \in \{50, 100, 200, 500, 1000, 1319\}$
without replacement, with $B{=}1000$ bootstrap resamples at each $M'$.

\begin{table}[htbp]
\centering
\small
\caption{Prompt-bootstrap CI half-width on $\hat{\mathcal R}_K$ shrinks as $1/\sqrt{M'}$ as the prompt sub-sample grows (headline cell: tulu3-8b-rlvr / gsm8k, $K{=}64$, $B{=}1000$). Last column is the half-width $\times\sqrt{M'}$; a near-constant ratio confirms the rate.}
\label{tab:appendix_M_convergence}
\begin{tabular}{rrr}
\toprule
$M'$ & CI half-width & $\sqrt{M'}\cdot$half-width \\
\midrule
50 & 0.0795 & 0.5624 \\
100 & 0.0347 & 0.3469 \\
200 & 0.0291 & 0.4116 \\
500 & 0.0178 & 0.3973 \\
1000 & 0.0149 & 0.4696 \\
1319 & 0.0129 & 0.4698 \\
\bottomrule
\end{tabular}
\end{table}

\subsection{Linking to Theorem~1}
\label{app:calibration:thm1}

The three calibrations above estimate the three terms of Theorem~1
on a per-cell basis:
\begin{enumerate}
\item Truncation in $K$: \cref{tab:appendix_saturation}, last-column gain
      $\hat{\mathcal R}_K(K_{\max}) - \hat{\mathcal R}_K(K_{\max}/2)$ upper-bounds the
      $K$-truncation residual under monotonicity.
\item Verifier-call variance in $L$: \cref{tab:L-sensitivity}, std
      column.
\item Prompt-sampling variance in $M$: \cref{tab:appendix_M_convergence}, CI
      half-width.
\end{enumerate}
On the headline GSM8K / Tülu-3-RLVR cell at $(M, K, L) = (1319, 64, 1)$
all three terms are simultaneously below \PH{$\epsilon$, e.g.\ 0.02},
which is the operating accuracy we report alongside the main-text
tables.

% =====================================================================
\section{Evaluator details}
\label{app:evaluator}

\subsection{Verifier-based evaluation (math and code)}

\paragraph{GSM8K.}
\texttt{src/verification/gsm8k.py} extracts the predicted number with a
four-pattern fallback chain applied in order:
\begin{enumerate}
\item \texttt{\#\#\#\# N} (the canonical GSM8K format),
\item \texttt{\textbackslash boxed\{N\}},
\item \texttt{"the answer is N"} (case-insensitive),
\item the last bare number anywhere in the text.
\end{enumerate}
The rightmost match within each pattern wins, since models routinely
restate intermediate numbers before committing. Both predicted and
ground-truth strings are run through the same extractor, so the verifier
is robust to a ground-truth distributed in any of the four formats.
The numeric token regex is the union of two alternatives — comma-grouped
(\texttt{1,234.56}) and bare (\texttt{1234.56}) — with optional leading
\texttt{\$} and optional trailing \texttt{\%}. Comparison is
\texttt{math.isclose(rel\_tol=1e-9, abs\_tol=1e-9)} so that
\texttt{72}, \texttt{72.0}, and \texttt{72.00} all match.

\paragraph{MATH (algebra).}
We use \texttt{math-verify} \texttt{>=0.9.0} for symbolic equivalence.
\texttt{src/verification/math.py} performs its own brace-balanced extractor
on the last \texttt{\textbackslash boxed\{...\}} in the generation
(handling nested braces, escape sequences \texttt{\textbackslash\{},
\texttt{\textbackslash\}}, \texttt{\textbackslash\textbackslash}, and
fall-back to earlier balanced openings if the rightmost open brace is
unmatched). We re-wrap both predicted and ground-truth expressions in
\texttt{\textbackslash boxed\{\ldots\}} before passing to
\texttt{math\_verify.parse} because the parser silently truncates bare
LaTeX like \texttt{2\textbackslash sqrt\{3\}} to its leading integer
\texttt{2}.

Parse and verify timeouts are 5\,s each; any timeout, parse error, or
internal exception in \texttt{math\_verify} resolves to $U=0$ (the
conservative fail-closed direction). \cref{tab:math-failures}
re-verifies every cached MATH generation for the three headline cells
($M{=}1000$ prompts, $K{=}64$ samples, $64{,}000$ pairs per cell) using
the instrumented verifier
\texttt{src.verification.math.verify\_with\_reason}, and tags each
$(x_i, y_{i,k})$ pair with its terminal state.

\begin{table*}[t]
\centering
\small
\setlength{\tabcolsep}{4pt}
\caption{MATH verifier failure-mode rates on the three headline cells,
percentage of $(x_i, y_{i,k})$ pairs. \texttt{correct} and
\texttt{incorrect} are the two semantic outcomes; \texttt{no
\textbackslash boxed} and \texttt{empty GT} are extraction failures
(failed before \texttt{math-verify} was called); \texttt{parse error},
\texttt{parse empty}, and \texttt{verify exception} are failures of the
\texttt{math-verify} library itself. All five non-semantic columns
resolve to $U{=}0$ (fail-closed).}
\label{tab:math-failures}
\begin{tabular}{l ccccccc r}
\toprule
Model & \shortstack[c]{no \\ \texttt{\textbackslash boxed\{\}}} & \shortstack[c]{empty \\ GT} & \shortstack[c]{parse \\ error} & \shortstack[c]{parse \\ empty} & \shortstack[c]{verify \\ exception} & incorrect & correct & $M{\cdot}K$ \\
\midrule
Tülu-3-8B-RLVR        &  4.9 & \textemdash & \textemdash & $<$0.1 & \textemdash & 29.3 & 65.8 & 64000 \\
Qwen2.5-7B-Instruct   &  2.6 & \textemdash & \textemdash & $<$0.1 & \textemdash &  7.7 & 89.8 & 64000 \\
Llama-3.1-8B-Instruct & 27.7 & \textemdash & \textemdash & $<$0.1 & \textemdash & 25.2 & 47.0 & 64000 \\
\bottomrule
\end{tabular}
\end{table*}

\paragraph{What the table says.}
The verifier's failure budget is overwhelmingly dominated by one mode:
the model never wrote a \texttt{\textbackslash boxed\{\ldots\}} commit
at all. Llama-3.1-8B-Instruct is the most affected at $27.7\%$ of
generations; Tülu-3-RLVR (instruction-tuned and RLVR'd to follow the
boxed-answer template) drops to $4.9\%$; Qwen2.5-7B (which natively
uses the template) is at $2.6\%$. The two \texttt{math-verify}-internal
failure modes that could in principle mask a correct answer — parser
exceptions and verifier exceptions — fire essentially never
($\leq 0.05\%$ combined across all three cells; the entire column is
absorbed by \texttt{parse empty}, which means the parser ran cleanly
and returned an empty expression list, typically on a non-mathematical
boxed value like \texttt{\textbackslash boxed\{?\}} or
\texttt{\textbackslash boxed\{\}}). The \emph{empty GT} column is a
defensive guard against a malformed ground truth and does not fire on
the algebra split.

\paragraph{Implication for the headline numbers.}
The fail-closed convention is therefore costing us at most a few tenths
of a percent of $U$ on these cells — the dominant cost is upstream
(the model didn't follow the answer format), not in the SymPy
back-end. In particular, the Llama-3.1-8B MATH $\bar U_K$ figure
underestimates the model's mathematical ability by roughly the size of
its $27.7\%$ no-\texttt{\textbackslash boxed} rate; the rational
expected utility $U^\circ_K$ is unaffected at $K{=}64$ because a single
correctly-formatted sample lifts the per-prompt indicator to $1$. This
asymmetric impact on $\bar U_K$ versus $U^\circ_K$ is itself part of
what $\hat{\mathcal R}_K$ measures — it shows up as inflated $\hat{\mathcal R}_K$
on Llama / MATH, exactly as the main-text table reports.

\paragraph{HumanEval.}
\texttt{src/verification/humaneval.py} concatenates the model's code
completion with the canonical \texttt{check} program plus a final
\texttt{check(<entry\_point>)} call, then executes the combined source
in a sandboxed subprocess with:
\begin{itemize}
\item a 5 s wall-clock timeout (\texttt{subprocess.run(timeout=5)});
\item a fresh \texttt{tempfile.TemporaryDirectory()} as cwd, so the
      candidate cannot read or pollute the repo or outputs tree;
\item a minimal environment that forwards only \texttt{PATH};
\item \texttt{start\_new\_session=True}, so SIGINT in the parent does
      not propagate and a fork bomb can be cleanly torn down.
\end{itemize}
$U{=}1$ iff the subprocess exits 0 within the timeout. Empty ground
truth fails closed at $U{=}0$ (an empty test program would otherwise
exit 0 vacuously). All 164 HumanEval problems have the full canonical
test suite of 7--10 assertions; the audit log
\texttt{outputs/logs/verifier/humaneval\_log.jsonl} records pass/fail
per assertion when the candidate raises before completing, allowing
post-hoc inspection of how much of the test suite a partial completion
satisfied. The aggregate unit-test coverage across HumanEval is the
canonical $\geq 7.7$ assertions/problem reported by \citet{chen2021codex}.

\paragraph{LiveCodeBench.}
\texttt{src/verification/livecodebench.py} parses the per-row test
cases from the (public + private) JSON blobs into a list of
\texttt{(testtype, input, output)} triples. For \texttt{stdin} tests
the candidate is invoked as a top-level script with \texttt{input}
piped to stdin and stdout compared (modulo trailing whitespace) to
\texttt{output}; for \texttt{functional} tests we instantiate
\texttt{Solution()} and call the method on the parsed arguments. Each
test has a 10 s per-test timeout; the whole problem has a 60 s
total-timeout cap. $U{=}1$ iff every test in the row passes. Code is
\emph{never} \texttt{exec()}'d in the parent process.

\paragraph{MathArena (AIME 2025 + BRUMO 2025).}
\texttt{src/verification/matharena.py} uses the same balanced-brace
extractor as the MATH verifier and the same
\texttt{math\_verify} equivalence check, with the following
defence-in-depth: if \texttt{math\_verify} fails to parse either side
(empty expression list or exception), we fall back to a
whitespace-insensitive string equality check
(\texttt{pred\_str.replace(" ", "") == gt\_str.replace(" ", "")}).
This is conservative — string equality is stricter than symbolic
equivalence, so the fallback can only reduce $U$, never inflate it.

\subsection{LLM-as-judge evaluation (UltraFeedback)}

\paragraph{Judge identity.}
Per the methodology memo and \texttt{scripts/run\_h1.py}
(lines 326--367) / \texttt{scripts/run\_h2.py} (lines 321--367):
\begin{itemize}
\item \textbf{H1 preference (strict-self):} the judge for each
      (model, prompt, candidate) triple is \emph{the same model} that
      generated the candidate. The H1 panel covers three judges:
      Tülu-3-8B-RLVR, Qwen2.5-7B-Instruct, Llama-3.1-8B-Instruct.
\item \textbf{H2 preference (fixed judge):} the judge is \emph{always}
      Tülu-3-8B-RLVR (\texttt{allenai/Llama-3.1-Tulu-3-8B}), regardless
      of which trajectory stage produced the candidate. This is a
      deliberate decision — judging SFT and DPO outputs through their
      own stages produces unreliable A/B parses (SFT is verbose,
      occasionally refuses), and we want the H2 claim to be ``the same
      judge's opinion across the trajectory''.
\end{itemize}
We do \emph{not} use GPT-4, Claude, or any other proprietary judge.
This is a choice driven by reproducibility — the entire H1/H2
preference panel can be re-run on a single 80 GB GPU with the model
weights named above and no API access.

\paragraph{Judge prompt.}
Reproduced verbatim in \cref{app:config:prompts}. The judge generates
at most 8 tokens at $\tau{=}0.7$; only the first non-whitespace
character is parsed (\texttt{A}/\texttt{B}/\texttt{T}).

\paragraph{Aggregation.}
With $L{=}5$ verdicts per (prompt, candidate) pair, the strict-majority
threshold is $\lceil L/2 \rceil = 3$: $U = 1$ if $\geq 3$ verdicts pick
the candidate, $U=0$ if $\geq 3$ pick the reference, $U=0.5$ otherwise
(no class clears the threshold).

\paragraph{Position-bias control.}
For each of the $L$ judge calls, the position of the candidate
(A or B) is decided by a fresh coin flip seeded deterministically from
\texttt{(seed, i, k, l)}, so the assignment is fully reproducible from
the cell seed alone (see
\texttt{src/verification/self\_judge.py}, lines 190--194).
\cref{tab:position-bias}
quantifies the residual position bias: the fraction of the $L=5$
verdicts that selected position A, marginalised across all
(prompt, candidate) pairs. A judge with no position bias would give
$\approx 0.5$.

\begin{table}[htbp]
\centering
\small
\caption{Residual position bias for the three H1 self-judges. ``A-pick
rate'' is the fraction of judge calls (across $M\cdot K\cdot L$
verdicts) where the judge picked the response in position A. A no-bias
judge gives 0.5.}
\label{tab:position-bias}
\begin{tabular}{lc}
\toprule
Judge model & A-pick rate \\
\midrule
Tülu-3-8B-RLVR (H1, H2)            & \PH{} \\
Qwen2.5-7B-Instruct (H1 only)      & \PH{} \\
Llama-3.1-8B-Instruct (H1 only)    & \PH{} \\
\bottomrule
\end{tabular}
\end{table}

\paragraph{Inter-rater agreement.}
We did not commission a human inter-rater study; the experimental
design here is faithful self-judgment, and the headline H1 preference
claim is about \emph{a model's own judgment} of its samples, not about
agreement with humans. As a sanity check we compute Krippendorff's
$\alpha$ across the $L{=}5$ verdicts of each judge call (treating the
five votes as five raters), which is what the recorded
\texttt{raw\_verdicts} field of the audit log enables. The headline
agreement numbers are in \cref{tab:irr}.

\begin{table}[htbp]
\centering
\small
\caption{Per-judge inter-rater agreement across the $L{=}5$ i.i.d.\
verdicts on each (prompt, candidate) pair. ``$\alpha$'' is
Krippendorff's $\alpha$ on the ternary scale $\{0, 0.5, 1\}$; ``mean
$|$majority margin$|$'' is the mean across pairs of the
absolute vote difference between the winning and losing classes.}
\label{tab:irr}
\begin{tabular}{lcc}
\toprule
Judge model & $\alpha$ (ternary) & mean $|$majority margin$|$ \\
\midrule
Tülu-3-8B-RLVR                & \PH{} & \PH{} \\
Qwen2.5-7B-Instruct           & \PH{} & \PH{} \\
Llama-3.1-8B-Instruct         & \PH{} & \PH{} \\
\bottomrule
\end{tabular}
\end{table}

\subsection{Answer-extraction failure cases}

\paragraph{GSM8K.}
The four-pattern fallback chain in
\texttt{src.verification.gsm8k.extract\_answer} is applied in priority
order: only the first pattern that returns a parseable number fires per
generation, and ``last bare number'' is reached only when
\texttt{\#\#\#\#}, \texttt{\textbackslash boxed}, and ``the answer is''
all fail. Re-extracting every cached GSM8K generation for the three
headline cells with the instrumented
\texttt{extract\_answer\_with\_pattern} returns, for each
$(x_i, y_{i,k})$ pair, both which pattern fired and whether the
extracted number was correct. \cref{tab:extract-fail} reports the
firing rate (\emph{fires}) and the conditional correctness given the
pattern (\emph{acc.}), across $M\cdot K = 1319 \cdot 64 = 84{,}416$
pairs per cell. A high \emph{fires} on \texttt{bare} \emph{combined
with} a low \emph{acc.}\ would indicate the verifier is leaning on a
noisy fallback; in practice it does not.

\begin{table*}[t]
\centering
\small
\setlength{\tabcolsep}{4pt}
\caption{Per-pattern firing and conditional-correctness rates of the
GSM8K extractor on the three headline cells. Rows are mutually
exclusive within each cell — only the first matching pattern fires per
generation. \emph{fires} sums to $100\%$ across the five columns;
\emph{acc.}\ is $P(U{=}1 \mid \text{pattern})$. The \texttt{no number
parsed} column is the residual: generations emitting no parseable
numeric token anywhere (these score $U{=}0$ by construction, so the
\emph{acc.}\ entry is undefined).}
\label{tab:extract-fail}
\begin{tabular}{l cc cc cc cc cc r}
\toprule
\multirow{2}{*}{Model} & \multicolumn{2}{c}{\texttt{\#\#\#\# $N$}} & \multicolumn{2}{c}{\texttt{\textbackslash boxed\{N\}}} & \multicolumn{2}{c}{``the answer is $N$''} & \multicolumn{2}{c}{\shortstack[c]{last bare \\ number}} & \multicolumn{2}{c}{\shortstack[c]{no number \\ parsed}} & \multirow{2}{*}{$M{\cdot}K$} \\
\cmidrule(lr){2-3} \cmidrule(lr){4-5} \cmidrule(lr){6-7} \cmidrule(lr){8-9} \cmidrule(lr){10-11}
 & fires & acc. & fires & acc. & fires & acc. & fires & acc. & fires & acc. & \\
\midrule
Tülu-3-8B-RLVR        & 99.9 & 86.2 & $<$0.1 & 62.5 & \textemdash & \textemdash & $<$0.1 & 25.0 & \textemdash & \textemdash & 84416 \\
Qwen2.5-7B-Instruct   & 84.9 & 90.6 & 14.7   & 90.9 & \textemdash & \textemdash & 0.4   & 72.1 & \textemdash & \textemdash & 84416 \\
Llama-3.1-8B-Instruct &  1.6 & 74.9 & $<$0.1 & 78.3 & 22.0        & 83.4        & 76.3  & 76.5 & $<$0.1      & \textemdash & 84416 \\
\bottomrule
\end{tabular}
\end{table*}

\paragraph{What the table says.}
The three models populate three distinct rows of the priority chain.
Tülu-3-RLVR sticks to \texttt{\#\#\#\#} on $99.9\%$ of generations —
the RLVR fine-tuning has effectively trained the canonical format into
the policy. Qwen2.5-7B splits between \texttt{\#\#\#\#} ($84.9\%$) and
\texttt{\textbackslash boxed} ($14.7\%$); accuracy is essentially
identical across the two patterns ($90.6$ vs $90.9$), confirming the
extractor is format-agnostic. Llama-3.1-8B reaches \texttt{\#\#\#\#}
only $1.6\%$ of the time, falls through ``the answer is'' for $22\%$
of generations, and uses the bare-number fallback for $76.3\%$. The
critical diagnostic — does the bare fallback inflate or deflate
correctness? — answers itself: on Llama, the four pattern accuracies
$(74.9, 78.3, 83.4, 76.5)$ are within a $9$-point band, so the
extractor's choice of pattern is not introducing a systematic bias.
The fallback chain is doing what it should: extracting the number, not
inventing it.

\paragraph{Residual unparseable rate.}
Across all three cells the \texttt{no number parsed} column is at most
$<0.1\%$ (Llama; absent for the other two). The four-pattern chain is
therefore a near-perfect cover of the number-emission distribution;
the GSM8K verifier's $U{=}0$ count is dominated by \emph{wrong
numbers}, not \emph{missing numbers}.

\paragraph{MATH and MathArena.}
Failure modes were broken down in \cref{tab:math-failures}. The
dominant mode is the upstream no-\texttt{\textbackslash boxed} rate
(model never committed); the SymPy parser and verifier paths
themselves are essentially never the cause of a $U{=}0$
($\leq 0.05\%$). MathArena uses the same extractor and the same
\texttt{math\_verify} call, plus a whitespace-insensitive
string-equality fallback (described in \cref{app:evaluator}); the
fallback is conservative (strict equality $\subseteq$ symbolic
equivalence) so it can only lower $U$, never raise it.

% =====================================================================
\section{Per-dataset and per-difficulty breakdown}
\label{app:per-dataset}

The main-text Table reports $\hat{\mathcal R}_K$ aggregated across (model, dataset).
This section breaks each headline cell down by problem difficulty
within the dataset, so the reader can verify that the rational gap is
not driven by a single difficulty stratum (in particular, not by the
easiest tier where $U^\circ_K$ saturates trivially).

\subsection{Stratification scheme}

\paragraph{GSM8K — step count.}
We bucket each problem by the number of reasoning steps in the
canonical \texttt{answer} field, where a ``step'' is a line of the
ground-truth solution containing an arithmetic operation
(\texttt{+}, \texttt{-}, \texttt{*}, \texttt{/}, \texttt{=}). The
buckets are $\{1\text{--}2, 3\text{--}4, 5\text{--}6, 7+\}$ steps,
which produces an approximately balanced distribution over the 1319
test problems (\PH{$M$/bucket: $\approx 280 / 480 / 380 / 180$}).

\paragraph{MATH (algebra) — level 1--5.}
We use the dataset's native \texttt{level} field directly, which gives
five buckets (Level 1 = easiest, Level 5 = hardest). Bucket sizes on
the 1187-problem algebra split are roughly
\PH{Level 1: $\approx 150$ / Level 2: $\approx 250$ / Level 3:
$\approx 280$ / Level 4: $\approx 290$ / Level 5: $\approx 210$}.

\paragraph{HumanEval — task ``difficulty''.}
HumanEval has no native difficulty annotation. We bucket by ground-truth
solution token length (a crude proxy for complexity), into terciles:
\textbf{easy} ($<60$ tokens, $\approx 55$ problems), \textbf{medium}
(60--120 tokens, $\approx 55$ problems), \textbf{hard} ($>120$ tokens,
$\approx 54$ problems). We do not attempt to use docstring complexity
or branching factor; the token-length proxy is sufficient to show that
$\hat{\mathcal R}_K > 0$ across the spectrum and is not concentrated at the
short-solution end.

\paragraph{MathArena and LiveCodeBench.}
The deployment-experiment splits are small ($M{=}60$ and $M{=}75$
respectively) so we do not stratify further; the per-difficulty
breakdown for these datasets is the headline cell. We include them
unstratified in \cref{tab:appendix_per_difficulty} for completeness.

\subsection{Headline breakdown}

\begin{table}[h]
\centering
\small
\setlength{\tabcolsep}{5pt}
\caption{Per-difficulty $\hat{\mathcal R}_K$ at $K{=}64$. MATH is bucketed by the official \texttt{level} field (1=easiest, 5=hardest); HumanEval is bucketed by canonical-solution-length tercile.}
\label{tab:appendix_per_difficulty}
\begin{tabular}{lll rrr}
\toprule
Dataset & Bucket & $n$ & Tülu-3-8B-RLVR & Qwen2.5-7B-Instruct & Llama-3.1-8B-Instruct \\
\midrule
MATH & L1 & 113 & 0.066 & 0.009 & 0.202 \\
 & L2 & 170 & 0.174 & 0.039 & 0.339 \\
 & L3 & 226 & 0.253 & 0.049 & 0.449 \\
 & L4 & 234 & 0.320 & 0.088 & 0.557 \\
 & L5 & 257 & 0.545 & 0.218 & 0.648 \\
\midrule
HumanEval & short & 56 & 0.392 & 0.093 & 0.337 \\
 & medium & 53 & 0.463 & 0.212 & 0.443 \\
 & long & 55 & 0.495 & 0.249 & 0.472 \\
\bottomrule
\end{tabular}
\end{table}

\textbf{Interpretation.} The take-away is that $\hat{\mathcal R}_K$ is bounded
away from zero in \emph{every} bucket, including the hardest tier
(Level 5 MATH, 7+ step GSM8K, hardest-tercile HumanEval), and is
not concentrated at any single end of the difficulty spectrum. The
saturation effect in $U^\circ_K$ on the easiest buckets is expected —
$U^\circ_K$ approaches the ceiling of 1 — but $\bar U_K$ saturates faster
there, so the gap shrinks at the easy end and \emph{grows} on
mid-difficulty problems where the reachable region remains rich but
sampling concentration drops. This pattern is consistent across all
three models.

% =====================================================================
\section{Failure-case analysis}
\label{app:failures}

This appendix illustrates the rational gap qualitatively: prompts where
$K{=}64$ contains at least one correct sample
($U^\circ_K(x_i) = 1$) but the model's sampling distribution
\emph{concentrates on incorrect answers} ($\bar U_K(x_i)$ small).
The candidate set is filtered by
$U^\circ_K(x_i)=1 \land \bar U_K(x_i) \le 0.1$ on the three headline
H1 cells. Per-prompt arrays come from the cell result JSONs
(\texttt{per\_prompt\_U\_circ\_K}, \texttt{per\_prompt\_U\_bar\_K});
the underlying generations come from the sample cache.

\paragraph{Spike-at-0 candidate prompts for the failure-case appendix.} The full per-prompt arrays in the result JSONs let any reader filter the same candidate set themselves; the rule is:
$U^\circ_K(x_i)=1$ \emph{and} $\bar U_K(x_i) \le 0.1$, i.e.\ the model can reach the correct answer in 64 samples but only on a small fraction of them. The candidate counts below come straight from the per-prompt arrays:
\begin{itemize}
  \item Tülu-3-8B-RLVR / GSM8K: 35 candidates out of 1319 prompts.
  \item Tülu-3-8B-RLVR / MATH: 68 candidates out of 1000 prompts.
  \item Tülu-3-8B-RLVR / HumanEval: 29 candidates out of 164 prompts.
  \item Qwen2.5-7B-Instruct / GSM8K: 22 candidates out of 1319 prompts.
  \item Qwen2.5-7B-Instruct / MATH: 15 candidates out of 1000 prompts.
  \item Qwen2.5-7B-Instruct / HumanEval: 9 candidates out of 164 prompts.
  \item Llama-3.1-8B-Instruct / GSM8K: 45 candidates out of 1319 prompts.
  \item Llama-3.1-8B-Instruct / MATH: 115 candidates out of 1000 prompts.
  \item Llama-3.1-8B-Instruct / HumanEval: 23 candidates out of 164 prompts.
\end{itemize}

From these candidates we select 4--5 traces per dataset and reproduce
(i)~the prompt, (ii)~the lowest-$k$ correct sample, and
(iii)~a representative incorrect sample (highest-$k$, $U=0$). The
selection step is editorial and not auto-generated; the scripted half
is the candidate filter shown above.
